\documentclass[tikz,border=4pt]{standalone}
\usepackage{ctex}
\usepackage{amsmath}
\usepackage{pgfplots}
\pgfplotsset{compat=1.18}

% p8-06b: 引入题——y=(k-2)x+2k 在不同 k 值下的图象，
% 特别突出 k=10/3 时图象经过 (1, 8) 的那一条
\begin{document}
\begin{tikzpicture}
\begin{axis}[
  width=9cm, height=8cm,
  axis lines=center,
  xlabel={$x$}, ylabel={$y$},
  every axis x label/.style={at={(axis cs:4.5,0)}, anchor=west},
  every axis y label/.style={at={(axis cs:0,9.5)}, anchor=south},
  xmin=-3, xmax=4.5,
  ymin=-3, ymax=9.5,
  xtick={-2,-1,1,2,3,4},
  ytick={-2,2,4,6,8},
  tick style={color=black},
  ticklabel style={font=\small},
  clip=true,
]
  % k=0: y=(-2)x+0 = -2x（斜率负，过原点）
  \addplot[gray!70, thick, domain=-2.5:3.5, samples=20] {-2*x};
  \node[gray, font=\small] at (axis cs:3.5, -6.2) {};  % out of view
  \node[gray, font=\small] at (axis cs:-1.5, 3.5) {$k=0$};

  % k=1: y=(-1)x+2 = -x+2（斜率负、截距正）
  \addplot[orange!80, thick, domain=-2.5:4, samples=20] {-x+2};
  \node[orange!80, font=\small] at (axis cs:4, -1.8) {$k=1$};

  % k=2 退化：y=0·x+4=4（水平线，常函数）
  \addplot[purple!60, thick, dashed, domain=-2.5:4, samples=20] {4};
  \node[purple, font=\small] at (axis cs:-2.3, 4.4) {$k=2$（退化）};

  % k=10/3: y=(10/3-2)x + 20/3 = (4/3)x + 20/3，经过 (1,8) ✓
  % 验证：(4/3)(1) + 20/3 = 4/3 + 20/3 = 24/3 = 8 ✓
  \addplot[red, very thick, domain=-3:3, samples=20] {(4/3)*x + 20/3};
  \node[red, font=\small\bfseries] at (axis cs:-2.6, 1.4) {$k=\frac{10}{3}$};

  % 特殊点 (1, 8) 高亮
  \addplot[only marks, mark=*, mark size=2.5pt, red] coordinates {(1, 8)};
  \node[red, font=\small, anchor=south west] at (axis cs:1.05, 8.1) {$(1, 8)$};

  % 虚线辅助 (1,8) 到坐标轴
  \draw[red!50, dashed, thin] (axis cs:1, 0) -- (axis cs:1, 8);
  \draw[red!50, dashed, thin] (axis cs:0, 8) -- (axis cs:1, 8);
\end{axis}
\end{tikzpicture}
\end{document}
